\documentclass[10pt,a4paper]{article}
\usepackage{amsmath,amssymb,amsthm,graphicx}
% Minimal retained preamble for S3 SIGMA article cards taken from
% sources/sigma-2023-028-cyt-and-skt-metrics-on-compact-semi-simple-l/pinned-original-source.tex
% (doi:10.3842/SIGMA.2023.028, SIGMA 19 (2023), CC BY 4.0).
%
% Omitted and why: the publisher \documentclass{sigma} block and the
% frontmatter commands - the sigma class is unavailable in the toolchain
% and the assembler supplies the article class.
% The theorem declarations, equation numbering and paper macros below
% are copied verbatim from the pinned source preamble.
% hyperref is added beyond the source preamble because the author's
% bibliography entries use \href and the unavailable sigma class is
% what provides it in the published version.
\usepackage{amsmath,amssymb,amsthm}
\usepackage{mathtools}
\usepackage{hyperref}
\allowdisplaybreaks
\numberwithin{equation}{section}
\newtheorem{Theorem}{Theorem}[section]
\newtheorem{Lemma}[Theorem]{Lemma}
\newtheorem{Proposition}[Theorem]{Proposition}
\newtheorem{Corollary}[Theorem]{Corollary}

\theoremstyle{definition}
\newtheorem{Definition}[Theorem]{Definition}
\newtheorem{Remark}[Theorem]{Remark}
\newtheorem{Example}[Theorem]{Example}

% The pinned source contains one stray Unicode non-breaking
% space U+00A0 (source line 451, inside the delivered proof
% span). XeLaTeX cannot typeset that character in cmr10; the
% two lines below make it an active character expanding to an
% ordinary space, which is the character the author's NBSP
% stands for. No source byte is modified and the delivered
% span remains a verbatim copy of the pinned source lines.
\catcode` =13
\let  \space
\title{CYT and SKT systems for invariant metrics; SKT+CYT implies Bismut flat (Corollary 5.1)}
\author{Fino, Anna, Grantcharov, Gueo}
\date{doi:10.3842/SIGMA.2023.028; CC BY 4.0}
\begin{document}
\maketitle
\setcounter{section}{5}
\setcounter{Theorem}{0}
\noindent\textit{Editorial note.} All delivered mathematical text is the author's own, copied from the pinned source file sources/sigma-2023-028-cyt-and-skt-metrics-on-compact-semi-simple-l/pinned-original-source.tex; no statement or proof step was rewritten, reordered or removed.\par\medskip
\noindent\textit{Editorial note.} The counted span is the author Corollary environment coroBismutflat with its own label and its own single proof environment (source lines 448-540). The author's phrase 'the last part of the theorem' at source line 476 (for the corollary's part (iii)) is copied verbatim.\par\medskip
\noindent\textit{Editorial note.} The Lambda definition (source lines 377-379) is delivered as a context span before the statement; the two sibling statements (Theorem 4.2, Corollary 3.4) are delivered as context spans after the statement so \textbackslash{}ref\{th4.1\} and \textbackslash{}ref\{th4.1bis\} resolve and the card keeps 5.1 (both print locally renumbered).\par\medskip
\noindent\textit{Editorial note.} No bibliography fragment is needed: the proof span carries no \textbackslash{}cite commands.\par\medskip
\noindent\textit{Editorial note.} The pinned source's stray Unicode non-breaking space U+00A0 (source line 451, inside the delivered proof span) is mapped to an ordinary space in the card preamble via newunicodechar; the delivered span text itself remains a verbatim line copy of the pinned source and no source byte is modified.\par\medskip
\noindent\textit{Editorial note.} Printed number anchored to the published PDF (SIGMA 19 (2023) 028, read from the stream-extracted text of sigma23-028.pdf downloaded from sigma-journal.com). The pinned tex has five \textbackslash{}section commands and no \textbackslash{}appendix and its section structure matches the published PDF one-to-one, so the unit 010 pinned-versus-published realignment does not trigger.\par\medskip
\noindent\textit{Editorial note.} Equation and theorem numbers inside the card are locally renumbered; every cross-reference is a symbolic \textbackslash{}eqref/\textbackslash{}ref and resolves inside the card.\par\medskip
\noindent\textit{Editorial scope.} Single environment card cut from the SIGMA article doi:10.3842/SIGMA.2023.028 (sigma-2023-028-cyt-and-skt-metrics-on-compact-semi-simple-l); statement and proof are the author's own text at source lines 430-445 and 448-540.\par\medskip
\noindent\textit{Allowed standard prerequisites (editorial).}\begin{itemize}
\item CYT characterization (Theorem 4.2). Delivered in context (source lines 420-424): g is CYT iff sum (d alpha - d(alpha circle Lambda)/lambda\_alpha) = 0.. Load-bearing: part (i) translates the CYT condition into the system (CYTsystem) via the simple-root expansion of alpha.
\item SKT criterion (Corollary 3.4). Delivered in context (source lines 214-220): the SKT criterion in root coordinates.. Load-bearing: part (ii) translates the SKT condition into the linear system (SKTsystem).
\item Simple-root data. Standard facts used by the author in-span: d alpha\_j = sum B(H\_alpha\_j, H\_alpha) e wedge Ie (source lines 466-469), B(H, H\_alpha) = alpha(H) (source lines 513-519), and the r x p simple-root coefficient matrix (x\_alpha\textasciicircum{}i) has rank r (source lines 534-538).. Load-bearing: the expansion of part (ii) and the invertibility argument of part (iii).
\end{itemize}\medskip
Let $g$ be a left-invariant $I$-Hermitian metric on $G$ which is also invariant by the right $T$-action. Then we can write $g (\cdot, \cdot ) = -B (\Lambda (\cdot), \cdot)$, where
$\Lambda$ is a positive-definite Hermitian matrix on $\mathfrak t$ and diagonal on $\sum_{\alpha} \mathfrak g_{\alpha}$.
In particular, $g (E_{\alpha}, E_{-\alpha}) = \lambda_{\alpha}$ and $g(E_{\alpha},E_{\beta}) =0$, for every $\alpha \neq - \beta$.

\begin{Corollary} \label{coroBismutflat}
Let $G$ be a compact semi-simple Lie group with Samelson's complex structure~$I$ determined by the set $P$ of positive roots associated to a maximal torus $T$ of $G$. Let $g$ be a~left-invariant $I$-Hermitian metric on $G$, which is also invariant under the right action of the torus~$T$ in the Tits fibration $G\rightarrow G/T$. Let $\alpha_1, \dots,\alpha_r$ be the choice of simple roots and $p$ be the cardinality of $P$.
Then
\begin{enumerate}
 \item[$(i)$] $g$ is CYT if and only if for every $\alpha \in P$ there exist $p$ positive real numbers $\lambda_{\alpha} > 0$, $\alpha \in P$, satisfying the $r$ equations
\begin{equation}\label{CYTsystem}
\sum_{\alpha \in P} \Bigg( x_\alpha^i - \sum_{j=1}^r \frac{x_\alpha^j}{\lambda_{\alpha}} \Lambda^i_j \Bigg) =0, \qquad i = 1, \dots, r,
\end{equation}
where $x_{\alpha}^j$ denotes the coefficient of $\alpha$ with respect to the simple root $\alpha_j$ and $\Lambda_i^j$ are the entries of $\Lambda$ in the basis $(\alpha_1, \dots,\alpha_r)$.
 \item[$(ii)$] $g$ is SKT if and only if there exist $r$ real numbers $ \tilde N_j$ and $p$ positive real numbers $ \tilde \lambda_{\alpha} >0$, $\alpha \in P$, satisfying the $p$ equations
\begin{equation}\label{SKTsystem}
\tilde \lambda_{\alpha} - 1 - \sum_{j = 1}^r \tilde N_j B(H_{\alpha_j}, H_{\alpha})=0, \qquad \alpha \in P.
\end{equation}
 \item[$(iii)$] If $g$ is both SKT and CYT, then $g$ is Bismut flat.
 \end{enumerate}
\end{Corollary}

\begin{Theorem} \label{th4.1}
Let $G$ be a compact semi-simple Lie group with a Samelson's complex structure determined by the set $P$ of positive roots associated to a maximal torus $T$ of $G$. Let $g$ be a~left-invariant $I$-Hermitian metric on $G$, which is also invariant under the right action of the torus $T$ in the Tits fibration $G\rightarrow G/T$, where $T$ is a maximal torus in $G$. We can write the fundamental form $F$ of the metric $g$
as $F = F|_{\mathfrak t} + \sum_{i=1}^m \sum_{\alpha \in P_i} \lambda_{\alpha} ^{(i)} e_{\alpha}\wedge I e_{\alpha}^{(i)}$, where $e_{\alpha}^{(i)}$, $Ie_{\alpha}^{(i)}$ are dual to~$X_{\alpha}$,~$Y_{\alpha}$ and $\lambda_{\alpha} ^{(i)}= g (E_{\alpha}, E_{-\alpha})$. Then
$g$ is CYT iff $\sum_{\alpha \in P} \big({\rm d}\alpha - \frac{1}{\lambda_{\alpha}} {\rm d}(\alpha\circ\Lambda) \big) = 0$, where $\Lambda$ is the positive operator such that $g |_{\mathfrak t} (\cdot, \cdot ) = -B (\Lambda (\cdot), \cdot)$.
\end{Theorem}

\begin{Corollary} \label{th4.1bis}
Let $G$ be a compact semi-simple Lie group with Samelson's complex structure $I$ as above. Let $g$ be a left-invariant $I$-Hermitian metric on $G$, which is also invariant under the right action of the torus $T$ in the Tits fibration $G\rightarrow G/T$, where $T$ is a maximal torus in $G$. We can write the fundamental form $F$ of the metric $g$
as $F = F|_{\mathfrak t} + \sum_{i=1}^m \sum_{\alpha \in P_i} \lambda_{\alpha} ^{(i)} e_{\alpha}^{(i)} \wedge I e_{\alpha}^{(i)}$, where $e_{\alpha}^{(i)}$, $Ie_{\alpha}^{(i)}$ are dual to $X_{\alpha}^{(i)}, Y_{\alpha}^{(i)}$ and $\lambda_{\alpha} ^{(i)}= g (E_{\alpha}^{(i)}, E_{-\alpha}^{(i)})$. Then
 $g$ is SKT iff $g |_{\mathfrak t} = - \sum_{i =1}^m \nu_i B|_{\mathfrak t_i}$ and $ \sum_{i =1}^m \sum_{\alpha \in P_i} \big(\lambda_{\alpha}^{(i)} - \nu_i\big) e_{\alpha}^{(i)} \wedge I e_{\alpha}^{(i)}$ is a linear combination of the basis $\omega_j = {\rm d} \alpha_j$ of $H^2(G/T, \mathbb R)$, where $\alpha_j$ are the simple roots.

Equivalently, $g$ is SKT iff ${\rm d}\big( \sum_{i =1}^m \sum_{\alpha \in P_i} \big(\lambda_{\alpha}^{(i)} - \nu_i\big) e_{\alpha}^{(i)}\wedge I e_{\alpha}^{(i)}\big)=0$ and $\Lambda |_{\mathfrak t} = \sum_{i=1}^m \nu_i {\rm Id}|_{\mathfrak t_i}$, for some constants $\nu_i$, $i= 1,2, \dots, m$.
\end{Corollary}

\begin{proof}
From Theorem \ref{th4.1}, it follows that $g$ is CYT if and only if there exist $p$ positive real number $\lambda_{\alpha}$, $\alpha \in P,$ such that
\[
\sum_{\alpha \in P}  \bigg ({\rm d}\alpha - \frac{1}{\lambda_{\alpha} } {\rm d}( \alpha \circ \Lambda ) \bigg) =0.
\]
Therefore, since $\alpha =\sum_{i = 1}^r x_\alpha^i \alpha_i$, with $x_\alpha^i$ positive integers and $\Lambda_i^j$ the entries of $\Lambda$,
 we have that the CYT condition is equivalent to
\[
\sum_{\alpha \in P} \Bigg( x_\alpha^i - \sum_{j=1}^r \frac{x_\alpha^j}{\lambda_{\alpha}} \Lambda^i_j \Bigg) {\rm d}\alpha_i = 0.
\]
Since the forms ${\rm d} \alpha_i$, $i= 1, \dots, r$, are linearly independent, we get the non-linear system \eqref{CYTsystem} of~$r$ equations in the $p$ variables $\lambda_{\alpha}$, $\alpha \in P$, and (i)~follows.

Let denote by $P_i$ the set of positive roots of the simple factor $G_i$ with respect to $\mathfrak t_i$.
To prove the second part of the corollary we use that by Corollary \ref{th4.1bis} the SKT condition is equivalent to the existence of $r$ real numbers $N_i$, $i = 1, \dots, r$ and $p$ positive numbers $\lambda_{\alpha}$, $\alpha \in P$, such that
\[
\sum_{i = 1}^m \sum_{\alpha \in P_i} \big(\lambda_{\alpha}^{(i)} - \nu_i\big) e_{\alpha}^{(i)} \wedge I e_{\alpha}^{(i)} = \sum_{i = 1}^r N_i \, {\rm d}\alpha_i
\]
and $\Lambda = \sum_{i =1}^m \nu_i {\rm Id} |_{\mathfrak t_i}$.
Since
\[
{\rm d}\alpha_j= \sum_{i = 1}^m \sum_{\alpha \in P_i} B\big(H_{\alpha_j}, H_{\alpha}^{(i)}\big) e_{\alpha} ^{(i)} \wedge I e_{\alpha}^{(i)}, \qquad j =1, \dots, r,
\]
we obtain
\[
\sum_{i =1}^m \sum_{\alpha \in P_i} \big(\lambda_{\alpha}^{(i)} - \nu_i\big) e_{\alpha}^{(i)} \wedge I e_{\alpha} ^{(i)}= \sum_{i =1}^m \sum_{j = 1}^r N_j \sum_{\alpha \in P_i} B\big(H_{\alpha_j}, H_{\alpha}^{(i)}\big) e_{\alpha}^{(i)} \wedge I e_{\alpha}^{(i)}.
\]
Therefore, using that the previous equality holds on every simple factor $\mathfrak g_i$ and rescaling $\lambda_{\alpha}^{(i)}$ and $N_i$ by $\nu_i$ we get the linear system \eqref{SKTsystem} of $p$ equations in the $p + r$ variables $\tilde N_i$, $i = 1, \dots, r,$ and $ \tilde \lambda_{\alpha}$, $\alpha \in P.$

Now to prove the last part of the theorem suppose that $g$ is both CYT and SKT. Then using that $\Lambda = \sum_{i =1}^m \nu_i {\rm Id} |_{\mathfrak t_i}$ and rescaling $\lambda_{\alpha}$ in \eqref{CYTsystem} in the same way as for the SKT case we have that the equations for the CYT condition become
\[
\sum_{\alpha \in P} x_\alpha^i\Bigg(1 - \sum_{j =1}^r \frac{1}{\tilde \lambda_{\alpha}} \Bigg) = 0, \qquad i = 1, \dots ,r.
\]
Let
\[
P = \{ \alpha_1, \dots, \alpha_r, \alpha_{r + 1}, \dots, \alpha_p \},
\]
where $\alpha_1, \dots, \alpha_r$ are the simple roots.
Imposing the SKT and CYT condition we get the following non-linear system of $p + r$ equations in the variables $\tilde \lambda_{j}$, $j =1, \dots, p$, $\tilde N_i$, $i = 1, \dots, r$,
\begin{equation}\label{newsystem}
 \begin{cases}
\displaystyle \tilde \lambda_{j}-1 - \sum_{i = 1}^r \tilde N_i B(H_{\alpha_i}, H_{\alpha_j})=0, \qquad j = 1, \dots, p,\\
\displaystyle \sum_{j=1}^p x_{\alpha_j}^i \bigg(1 - \frac{1}{\tilde \lambda_{j}} \bigg) =0, \qquad i = 1, \dots, r.
\end{cases}
\end{equation}
Note that the solution $\tilde N_i =0$, $i=1, \dots, r$, $\tilde \lambda_j = 1$, $j = 1, \dots, p$, corresponds to a Bismut flat metric.
Since $ \tilde \lambda_j \neq 0$, for every $j = 1, \dots, p$, if we multiply
\[
\sum_{j=1}^p x_{\alpha_j}^i \bigg(1 - \frac{1}{\tilde \lambda_{j}} \bigg) =0, \qquad i = 1, \dots, r,
\]
by the product of all the $\tilde \lambda_j$,we get
\[
\sum_{j=1}^p x_{\alpha_j}^i \tilde \lambda_1 \cdots \hat{\tilde {\lambda_j}} \cdots \tilde \lambda _p \big(\tilde \lambda_j - 1\big) =0, \qquad i = 1, \dots, r,
\]
where $\hat{ \tilde {\lambda_{j} }}$ denotes that the element $\tilde \lambda_{j}$ is removed in the product. We know that
\[
\tilde \lambda_{j}-1 = \sum_{k = 1}^r \tilde N_k B(H_{\alpha_k}, H_{\alpha_j})
\]
for every $j = 1, \dots, p$. Therefore
\[
\sum_{j=1}^p x_{\alpha_j}^i \tilde \lambda_1 \cdots \hat{\tilde {\lambda_j}} \cdots \tilde \lambda _p \Bigg (\sum_{k = 1}^r \tilde N_k B(H_{\alpha_k}, H_{\alpha_j})\Bigg)=0, \qquad i = 1, \dots, r,
\]
 which can be viewed as a homogeneous system of $r$ equations in the $r$ variables $\tilde N_i$, $i = 1, \dots, r$. The $r \times r$ associated matrix $M = M\big(\tilde \lambda_1, \dots, \tilde \lambda_p\big)$ has $(i,k)$-th entry given by
\[
\sum_{j=1}^p x_{\alpha_j}^i \tilde \lambda_1 \cdots \hat{\tilde {\lambda_j}} \cdots \tilde \lambda _p B(H_{\alpha_k}, H_{\alpha_j}).
\]
We can rewrite $B(H_{\alpha_k}, H_{\alpha_j})$, $k =1, \dots, r$, $j = 1, \dots, p,$ in terms of $B(H_{\alpha_k}, H_{\alpha_j})$ with $k, j \in \{1, \dots, r \}$, using
 that $B (H, H_{\alpha}) = \alpha (H)$ for every $H \in \mathfrak t$. Indeed, if $\alpha = \sum_{l=1}^r x_{\alpha}^l \alpha_l$ we have,
\[
B (H_{\alpha_i}, H_{\alpha}) = \alpha (H_{\alpha_i}) =
 \sum_{l=1}^r x_{\alpha}^l \alpha_l (H_{\alpha_i}) = \sum_{l=1}^r x_{\alpha}^l B(H_{\alpha_i}, H_{\alpha_l}),
\]
 for every $i = 1, \dots, r$.
 As consequence the $(i,k)$-th entry of the matrix $M = M\big(\tilde \lambda_1, \dots, \tilde \lambda_p\big)$ can be written as
\begin{gather*}
\sum_{j=1}^p x_{\alpha_j}^i \tilde \lambda_1 \cdots \hat{\tilde {\lambda_j}} \cdots \tilde \lambda _p B(H_{\alpha_k}, H_{\alpha_j}) = \sum_{j=1}^p x_{\alpha_j}^i \tilde \lambda_1 \cdots \hat{\tilde {\lambda_j}} \cdots \tilde \lambda _p B(H_{\alpha_j}, H_{\alpha_k})\\
\hphantom{\sum_{j=1}^p x_{\alpha_j}^i \tilde \lambda_1 \cdots \hat{\tilde {\lambda_j}} \cdots \tilde \lambda _p B(H_{\alpha_k}, H_{\alpha_j})}{} = \sum_{j=1}^p x_{\alpha_j}^i \tilde \lambda_1 \cdots \hat{\tilde {\lambda_j}} \cdots \tilde \lambda _p \sum_{l=1}^r x_{\alpha_j}^l B(H_{\alpha_l}, H_{\alpha_k})\\
 \hphantom{\sum_{j=1}^p x_{\alpha_j}^i \tilde \lambda_1 \cdots \hat{\tilde {\lambda_j}} \cdots \tilde \lambda _p B(H_{\alpha_k}, H_{\alpha_j})}{} = B\Bigg( \sum_{l=1}^r \sum_{j=1}^p x_{\alpha_j}^i \tilde \lambda_1 \cdots \hat{\tilde {\lambda_j}} \cdots \tilde \lambda _p x_{\alpha_j}^l H_{\alpha_l}, H_{\alpha_k}\Bigg),
\end{gather*}
where $i , k = 1, \dots, r$. Since
\[
 \sum_{j=1}^p x_{\alpha_j}^i \tilde \lambda_1 \cdots \hat{\tilde {\lambda_j}} \cdots \tilde \lambda _p x_{\alpha_j}^l = \tilde \lambda_1 \cdots \tilde \lambda_p \sum_{j=1}^p x_{\alpha_j}^i \frac{1} { \tilde \lambda_j} x_{\alpha_j}^l,
\]
 it will be sufficient to show, for every $ \tilde \lambda_j$, the invertibility of the matrix $L$, whose $(i,l)$-entry is given
\[
 \sum_{j=1}^p x_{\alpha_j}^i \tilde \lambda_1 \cdots \hat{\tilde {\lambda_j}} \cdots \tilde \lambda _p x_{\alpha_j}^l = \tilde \lambda_1 \cdots \tilde \lambda_p \sum_{j=1}^p x_{\alpha_j}^i \frac{1} {\tilde \lambda_j} x_{\alpha_j}^l.
\]
Since the $r \times p$ matrix $(x_{\alpha}^i)$ has rank $r$, the $r$ vectors
\[
v_1 = \big(x^1_{\alpha_1}, \dots, x^1_{\alpha_p}\big), \quad \dots, \quad v_r = \big(x^r_{\alpha_1}, \dots, x^r_{\alpha_p}\big),
\]
are linearly independent. Using that for every $j=1, \dots, p,$ $ \tilde \lambda_j > 0$, we get that $\sum_{j=1}^p \frac{1} { \tilde \lambda_j} x_{\alpha_j}^i x_{\alpha_j}^l $ can be viewed as a positive definite scalar product on $\mathbb R^p$ with weights $\frac{1}{ \tilde \lambda_j}$ of the two vectors~$v_i$ and~$v_l$. Therefore,
the matrix $L$ is invertible and $\tilde N_i =0$, for every $i = 1, \dots, r$. By \eqref{newsystem}, it follows that $\tilde \lambda_j = 1$, for every $j =1, \dots, p$.
\end{proof}

\end{document}
